% СПО ЛКС - Распредеденные очереди сообщений
% Мельников Максим (с) 2011
% Пынькин Д.А. (с) 2012
\documentclass[ignorenonframetext, hyperref={pdftex, unicode}]{beamer}

\usepackage[russian]{babel}
\usepackage[utf8]{inputenc}
\usepackage[T1]{fontenc}

\usepackage{beamerthemesplit}

%\usetheme{Pittsburgh}
%\usecolortheme{dolphin}

\usetheme{Copenhagen}
\usecolortheme{wolverine}


\usepackage{ulem}
%\usepackage{html}

\usepackage{verbatim}

\usepackage{ucs}
\usepackage{listings}
\lstloadlanguages{C, make, bash}

\lstset{escapechar=`,
	extendedchars=false,
	language=C, 
	tabsize=2, 
	columns=fullflexible, 
%	basicstyle=\scriptsize,
	keywordstyle=\color{blue}, 
	commentstyle=\itshape\color{brown},
%	identifierstyle=\ttfamily, 
	stringstyle=\mdseries\color{green}, 
	showstringspaces=false, 
	numbers=left, 
	numberstyle=\tiny, 
	breaklines=true, 
	inputencoding=utf8x,
	keepspaces=true,
	morekeywords={u\_short, u\_char, u\_long, in\_addr}
	}

\definecolor{darkgreen}{cmyk}{0.7, 0, 1, 0.5}

\lstdefinelanguage{diff}
{
    morekeywords={+, -},
    sensitive=false,
    morecomment=[l]{//},
    morecomment=[s]{/*}{*/},
    morecomment=[l][\color{darkgreen}]{+},
    morecomment=[l][\color{red}]{-},
    morestring=[b]",
}

\title[СПОЛКС (http://goo.gl/32cTB)]{Системное программное обеспечение\\локальных компьютерных сетей}
\author{Денис Пынькин}
\date{2011 -- 2012}
%\institution[БГУИР]{Белорусский государственный университет информатики и радиоэлектроники}
%\logo{\includegraphics[width=1cm]{logo-kafEVM.png}}


\subtitle{Очереди сообщений в распределённых системах}

\begin{document}

\mode<all>{\begin{frame}
\titlepage
\begin{center}
e-mail: denis.pynkin@bsuir.by\\
\end{center}
\begin{center}
{\bfseries http://goo.gl/32cTB}

{\tiny СЧАСТЬЕ ДЛЯ ВСЕХ, ДАРОМ, И ПУСТЬ НИКТО НЕ УЙДЕТ ОБИЖЕННЫЙ!\\
(c)Стругацкие, Пикник на обочине}
\end{center}
\end{frame}
}

%
% Далее начинается сама презентация
%
\frame{
    \begin{center}
    Выражается благодарность за подготовку материалов лекции:

    \huge
    Максиму 'max\_posedon' Мельникову

    \normalsize
    m\_melnikau@wargaming.net

    \end{center}
}


\frame{
    \frametitle{План}
    \tableofcontents
}

\section{Обзор различных MQ}
\frame{
    \frametitle{Очереди сообщений}
    \begin{block}{Что это такое?}
    Это архитектура и ПО промежуточного уровня, которое занимается сбором, хранением 
    и маршрутизацией (распределением) сообщений между компонентами. 
    \end{block}

    \begin{block}{Зачем?}
    Организация очереди сообщений помогает раз-балансировать нагрузку между различными узлами сети, 
    избавиться от единой точки отказа (SPOF),
    выполнять бизнес-логику приложения асинхронно, 
    повысить скорость ответа системы и многое другое.
    \end{block}
}

\frame{
    \frametitle{Apache ActiveMQ}
    \begin{itemize}
    \item Реализован на Java
    \item Полностью реализует Java Message Service 1.1 (JMS)
    \item "Enterprise Features": кластеризация, журнал операций, ...
    \end{itemize}
}

\frame{
    \frametitle{ZeroMQ}
    \begin{itemize}
    \item Реализован на C++
    \item Работает без сервера
    \item Низкоуровневое API
    \end{itemize}
}

\frame{
    \frametitle{PgQ}
    \begin{itemize}
    \item Реализован на C/SQL как расширение к PostgreSQL
    \item Тесная интеграция с базой данных (транзакции)
    \item Skype
    \end{itemize}
}

\frame{
    \frametitle{AMQP}
    \begin{block}{AMQP}
    Advanced Message Queueing Protocol   
    \end{block}

    \begin{block}{Apache Qpid}
    \begin{itemize}
    \item Реализация на C++
    \item RedHat
    \end{itemize}
    \end{block}

    \begin{block}{RabbitMQ server}
    \begin{itemize}
    \item Реализация на Erlang
    \item VMWare
    \end{itemize}
    \end{block}
}

\section{Архитектура AMQP}
\frame{
    \frametitle{Используемая терминология}
    \begin{itemize}
    \item Брокер (broker)
    \item Сообщение (mesasage)
    \item Точка обмена (exchange)
    \item Очередь (queue)
    \end{itemize}
}

\frame{
    \frametitle{Hello World}
    \begin{figure}[htb]
    \includegraphics[width=\textheight]{hello.png}
    \end{figure}
}

\frame{
    \frametitle{Work queues}
    \begin{figure}[htb]
    \includegraphics[width=\textwidth]{queues.png}
    \end{figure}
}

\frame{
    \frametitle{Publish/Subscribe}
    \begin{figure}[htb]
    \includegraphics[width=\textwidth]{pubsub.png}
    \end{figure}
}

\frame{
    \frametitle{Routing}
    \begin{figure}[htb]
    \includegraphics[width=\textwidth]{routing.png}
    \end{figure}
}

\frame{
    \frametitle{Topics}
    \begin{figure}[htb]
    \includegraphics[width=\textwidth]{topics.png}
    \end{figure}
}

\frame{
    \frametitle{RPC}
    \begin{figure}[htb]
    \includegraphics[width=\textwidth]{rpc.png}
    \end{figure}
}

\mode<all>{\begin{frame}{}
\Huge
\begin{center}
	Спасибо за внимание!
	\bigskip
	Вопросы?
\end{center}
\normalsize
\end{frame}
}

\end{document}
